\begin{recipe}{Drollenkoekjes}
\kicker[Bakken,Zoet,Kinderen]{Bakken · zoet · kinderen}
\index[register]{Drollenkoekjes}
\index[register]{Cacao}

\meta{1 uur koelen + 15--20 minuten bakken \dvd 12 stuks \dvd Oven 175\,°C}

\marginimage{images/drollenkoekjes}
\lettrine{C}{hocoladekoekjes} in de vorm van een drol, als kindertraktatie. Ze
zijn niet moeilijk te maken en vallen meestal goed bij kinderen. Het deeg moet
wel minstens een uur rusten in de koelkast.

\blockrule{Ingrediënten}
\begin{ingredients}
  \ing{360 g}{patentbloem}
  \ing{200 g}{ongezouten roomboter, kamertemperatuur (ca.\ 56\%)}
  \ing{150 g}{donkere basterdsuiker (ca.\ 42\%)}
  \ing{2 zakjes}{vanillesuiker}
  \ing{1}{ei}
  \ing{40 g}{cacao}\index[register]{Cacao}
  \ingb{eetbare oogjes}
  \ingb{poedersuiker, voor het glazuur}
\end{ingredients}

\blockrule{Bereiding}
\tip{Hoe kouder het deeg, hoe minder de koekjes uitlopen. Een nachtje in de
koelkast werkt beter dan een uurtje.}
\begin{steps}
  \item Mix de boter met de basterdsuiker en vanillesuiker tot een gladde massa.
  \item Voeg het ei, de bloem en cacao toe en kneed tot alles goed gemengd is.
  \item Maak een bol van het deeg, wikkel in vershoudfolie en laat minimaal 1 uur
        in de koelkast rusten.
  \item Verwarm de oven voor op 175\,°C (boven- en onderwarmte). Verdeel het deeg in 12 gelijke stukken.
  \item Rol elk stuk tot een lange slang met een puntig uiteinde. Draai de slang
        op tot een drolletje en zet het puntje iets omhoog.
  \item Bak 15--20 minuten in de oven. Laat volledig afkoelen.
  \item Meng een beetje poedersuiker met een paar druppels water tot een dik
        glazuur. Plak de oogjes ermee op de koekjes.
\end{steps}
\end{recipe}
